% Contenido del póster: Ran, Rau y Ziegler (2025), Management Science. DOI 10.1287/mnsc.2023.01538
% Elaborado con apoyo de IA a partir del resumen oficial del artículo (OpenAlex), 2026-11-12.
% El estudiante debe leer el texto completo, completar lo marcado y reescribir con sus palabras.
\def\DOI{10.1287/mnsc.2023.01538}
\def\TituloEs{A veces, siempre, nunca: claridad regulatoria y desarrollo del financiamiento digital}
\def\TituloOriginal{Sometimes, always, never: Regulatory clarity and the development of digital financing}
\def\Autores{Ran, Rau y Ziegler}
\def\Anio{2025}
\def\Revista{Management Science}

\def\Problema{En la última década, muchos países regularon el crowdfunding, pero con distinto grado de detalle: algunos dictaron normas explícitas, otros reglas generales y otros ninguna. La pregunta es si importa solo que exista una regulación o también qué tan clara y precisa es.}
\def\Objetivo{Examinar cómo el nivel de detalle de las normas legales influye en el volumen del crowdfunding de deuda y de capital.}
\def\Metodo{Los autores usan una muestra global de financiamiento digital y miden el detalle regulatorio con indicadores aproximados. Comparan tres grupos de países: los que tuvieron regulación solo en parte del periodo («a veces»), los que la tuvieron siempre y los que nunca la tuvieron, y estiman una serie de regresiones de diferencias en diferencias.}
\def\Resultados{\begin{itemize}
  \item Las regulaciones claras y explícitas aumentan de forma significativa el volumen del crowdfunding.
  \item La relación positiva aparece en el crowdfunding de deuda; en el de capital el efecto es pequeño.
  \item La claridad regulatoria favorece la creación de nuevas plataformas, en lugar de concentrar la actividad en las que ya existían.
\end{itemize}}
\def\Conclusion{No basta con regular: el grado de detalle de la norma importa. Las reglas explícitas se asocian con más crowdfunding de deuda y con la entrada de nuevas plataformas, lo que da a los reguladores una herramienta concreta para desarrollar este mercado.}
\def\Relevancia{El Perú reguló el financiamiento participativo financiero con el Decreto de Urgencia N.\textsuperscript{o}~013-2020 y el reglamento de la SMV de 2021, y lo modificó en noviembre de 2025. Según este estudio, cabría esperar más crédito por plataformas y más competidores. Con datos de la SMV, el monto anual pasó de S/~\cifra{monto2022} millones en 2022 a S/~\cifra{monto2025} millones en 2025, pero solo \cifra{plataformasConFinanciamiento} de \cifra{plataformasAutorizadas} sociedades administradoras financiaron proyectos y una concentra el \cifra{participacionLider}~\% del monto: el volumen creció, pero la entrada de plataformas no.}
\def\PalabrasClave{regulación, claridad regulatoria, crowdfunding de deuda, plataformas, diferencias en diferencias}
% --- Ficha del estudio (herramientas/posteres/ficha_articulos.py) ---
\def\Ambito{Muestra global de países}
\def\Enfoque{Empírico: diferencias en diferencias}
\def\FiguraPeru{1}
% --- fin de la ficha ---
